\scanpage{63}

\begin{center}- 149 -\end{center}
\begin{flushright}IX\end{flushright}

\uline{en associant à tout} $i\in I$ \uline{la section} $S_i$ \uline{de} $P$ \uline{sur} $S$ (\uline{nulle en} $\eta$) \uline{définie par le faisceau inversible} $\uline{O}_X(X_{\unsure{o|0}}^i)$ \uline{sur} $S$. (NB \uline{Comme les anneaux locaux de} $X$ \uline{sont factoriels en hypothèse}, \uline{les cycles} $X_{\unsure{o|0}}^i$ \uline{sont bien des diviseurs}, \uline{ce qui donne un sens à} $\uline{O}_X(X_i^{\unsure{o|0}})$ \uline{comme le faisceau inversible associé au diviseur en question} (EGA IV 21.2.8.1).) \uline{On a}
\[
(12.1.6)\qquad P^{\unsure{o|0}}\cap E=\text{\uline{section nulle de }}P\quad.
\]

c) \uline{Sous les conditions de} b), \uline{le quotient}
\[
(12.1.7)\qquad Q=P/E
\]
\uline{est représentable par un schéma en groupes lisse et séparé sur} $S$, \uline{et ce dernier est} néronien, \uline{i.e. satisfait la propriété universelle} (1.1.2) \uline{en termes de sa fibre spéciale}.

d) \uline{Sous les conditions de} c), \uline{supposons de plus que la fibre géométrique générique de} $X$ \uline{soit} intègre. \uline{Alors on a une suite exacte canonique}
\[
(12.1.8)\qquad
0\longrightarrow A\longrightarrow Q\xrightarrow{\deg}\mathbb Z_S\longrightarrow0\quad,
\]
\uline{où} $\deg$ \uline{est induit par l'homomorphisme} ``\uline{degré}'' $P\longrightarrow\mathbb Z_S$, \uline{et où} $A$ \uline{est un schéma en groupes lisse séparé et} de type fini \uline{sur} $S$, \uline{qui est également néronien i.e. qui est un modèle de Néron de}
\[
(12.1.9)\qquad A_\eta=P_\eta^{\unsure{o|0}}=\uline{\operatorname{Pic}}^{\unsure{o|0}}_{X_\eta/\eta}\quad.
\]
\uline{D'autre part, l'homomorphisme composé} $P^{\unsure{o|0}}\longrightarrow P\longrightarrow Q$ \uline{induit un isomorphisme}
\[
(12.1.10)\qquad P^{\unsure{o|0}}\xrightarrow{\sim}A^{\unsure{o|0}}\quad.
\]

\margin{467}
\margin{467}
\margin{62}
